% !TeX root = ../main.tex

\chapter{Design Methodology}
\label{ch:methodology}

The design process involved an iterative engineering method
starting with requirements determination and preliminary
design and moving on to detailed circuit design, prototyping,
implementation, and testing.

\section{Engineering Analysis}

Component selection, schematic drawing, PCB designing,
power budgeting, current sense resistor calculation,
and electrical operating limitations verification were
some of the tasks performed during the preliminary and
detailed designs.

The calculations based on Ohm's law, power dissipation,
ADC resolution, component tolerance values, and signal
isolation techniques were used to determine the circuit
requirements. Calculations are provided in
Appendix~\ref{app:design-calculations}.

\section{Design Trade-offs}

Among the main trade-offs of the design were resolution,
power dissipation, isolation, cost, and development time.

For measuring current, high precision shunt resistors were
selected balancing resolution against voltage and power
dissipation. Independent isolated channels were selected to
maintain the isolation/floating rails of the measured supplies although
this entailed extra circuitry.

For speeding up the development process, Arduino GIGA
platform with its set of peripherals and display capability
was used. Wired Ethernet was chosen for ease of integration into a control system.

\section{Broader Considerations}

In addition to that, safety of electricity use,
accuracy and reliability of measurements, maintainability,
cybersecurity, and responsible components sourcing were
among other considerations of the design which will be
addressed in Chapter~\ref{ch:impacts}.

Both preliminary and detailed design were then improved
through prototyping process discovering some hardware and
communication issues which required changes including
changes to the I2C circuit.